% 序盤（じょばん）：配牌から6巡目という一段目の打牌

\documentclass[../nankiru_main.tex]{subfiles}

% override main tex render option
\renewcommand{\renderOption}{1}

\begin{document}

\subfilepreamble{序盤戦略}

\subsection{バラバラの配牌の打点上昇}

バラバラの配牌が来て和了りにくく見えそうけど、稀に和了れた時には満貫以上の高い打点にしたい。打点があれば、仮に先制立直を打たれても一向聴から押していくこともあるかもしれない\citep[][pp. 40--47]{tetsusenryaku}。

\tehai{135m 127p 25899s 22 -4z}[東1局　西家　ドラ \mj{4z}]

和了れなさそうな配牌。チートイ、チャンタ、下の三色が見える。チャンタに向かう \mj{5m} や \mj{5s} を切ると \mj{0m} と \mj{0s} の受けが無くなり、\mj{7p} にくっつくと三色が消滅する。門前で行けば赤引きの可能性を残し \mj{7p} を切ってチートイ立直を目指す。ちなみに \mj{2s} はチャンタや三色のどちらにも使えるので、\mj{5s} より残したい。チャンタと三色は仕掛けると打点が下がるので要注意である。

\tehai{11224m 246p 389s 11 -4z}[東1局　西家　ドラ \mj{4z}]

面子手で和了ると辺張より浮き牌を切る。ですがこんな聴牌まで遠い階段なら東をスルーしてもあり、チートイに目指すなら \mj{8s} を先切りにする。

%--------------------------------------------
\subsection{高くてバラバラの配牌}

字牌が重なる可能性が低いが、重なると鳴きにより数牌で面子を整えるのより早いので、役牌に絡む役（チャンタとか）も見込む。もちろん、毎回字牌が重なるそして誰かが鳴かせるわけではない、それでも字牌が重なるチャンスを見逃さない。誰かが字牌を切っても安易に合わせもしない\citep[][pp. 56--63]{tetsusenryaku}。

\tehai{1116m 127p 2679s 16 -7z}[東1局　西家　ドラ \mj{1m}]

立直に目指せば浮き牌の \mj{6m} と \mj{2s} を比べると \mj{6m} の方がわずかな強い。チャンタに向かうなら \mj{3p} でもチーして \mj{6m6s} を落とす。

\subsubsection*{他家が早そうな局面}

\begin{rittai}[h]
  \centering
  
  \drawrittai{
    jicha/seat = 西,
    jicha/hand = 679m 69p 9s 1223446z -3m,
    jicha/river = 23^s 4p,
    %
    toimen/hand = XXXXXXXXXX -777'z,
    toimen/river = 1s 5^z 9^6s 5p,
    %
    kami/river = 45^z 2m 7p,
    %
    shimo/river = 3z 4m 8^p,
    %
    kyoku = 東4局,
    dora = 3m
  }
  
  \caption{手牌が多くて他家が早そうな局面}
  \label{rittai:joban-far}
\end{rittai}

混一色が見える手牌だが親が1副露、そして他家がもう中張牌切っている。横移動を期待するのが最善なので、Suphxは打 \mj{4z} とした\citep[][p.~35]{suphx}。

%--------------------------------------------
\subsection{辺張外し}

いい待ちになりやすければ、または打点が上昇すれば、序盤で辺張を外すのは正解。打点十分であれば、または辺張を落としても打点やアガリ率も上がらないなら、そのまま聴牌を取るべき。特に、巡目が深くなれば（6〜7巡以降）、辺張落とすのは不利になる。ただし、その場合、辺張待ちになったら立直しない方がおすすめ\citep[][pp. 48--55]{tetsusenryaku}。

\tehai{4678m 3589p 11067 -7s}[東1局　西家　3巡目　ドラ \mj{2m}]

タンヤオがないため、辺張を落としても打点はよほどよくならなさそう。聴牌のときの最終形の想定すれば、辺 \mj{7p} を避けるもっといい待ちを作る。

\tehai{4678m 3589p 11067 -7s}[東1局　西家　3巡目　ドラ \mj{1s}]

ドラ3で打点はもう十分である場合は、素直に聴牌に向かうべき。ですが、ソウズの部分が \mj{114067s} になったら伸びやすいので辺張を外すのはおすすめ。

%--------------------------------------------
\subsection{嵌張外し}

\tehai{344m 79m 5667p 2267 -8s}[東1局　西家　4巡目　ドラ \mj{2p}]

マジョリティは \mj{79m} の嵌張を落としてタンヤオを確定するが、\mj{8m} が先に埋まったメンピンの可能性が消滅する。\mj{79m} を落としても \mj{344m} からもう一つ面子ができる可能性が低いので、\mj{4m} 切りがおすすめ\citep[][pp. 92--97]{tetsusenryaku}。

%--------------------------------------------
\subsection{両面聴牌外し}

\tehai{34m 234p 3335567s 2 -2z}[東1局　西家　3巡目　ドラ \mj{7m}]

\tehai{34m 2345p 333567s 2 -2z}[東1局　西家　3巡目　ドラ \mj{7m}][2-4-1]

\tehai{34m 2346p 333567s 2 -2z}[東1局　西家　3巡目　ドラ \mj{7m}][2-4-2]

全ての状況は早めの3巡目。そのまま立直すれば打点はただの1300ですが、\mj{22z} を落としタンヤオや三色やピンフで打点は大きく上昇する。\ref{mj:2-4-1}は\mj{22z} を落として \mj{25m} を引けばまだ三色やピンフに替わり機会が多いのでダマはおすすめ。\ref{mj:2-4-2}が \mj{6p} 単騎になる可能性があるので、それを避けて即リーがおすすめ\citep[][pp. 68--73]{tetsusenryaku}。

%--------------------------------------------
\subsection{対子外し}

\tehai{123 89m 79p 155s 244 -4z}[南4局　南家　1巡目　ドラ \mj{1s}]

チャンタが見えるが、愚形ばかりな手牌。そのままで立直すれば大体2600に過ぎない。ですので、\mj{5s} の対子を落とし、\mj{2z} 重なりや \mj{4z} 槓で打点を最大限にするのはおすすめ\citep[][pp. 74--79]{tetsusenryaku}。

%--------------------------------------------
\subsection{アガリトップ}

\tehai{588m 12699p 266s 25 -7z}[南4局　西家　1巡目　ドラ \mj{5z}　アガリトップ]

一気に見ると和了れなさそうな配牌であるが、供託があるのでアガったらトップになる、ワンスリー\footnote{トップ目のポイントは$+30$にして、二着目のは$+10$、三着目のは$-10$、ラス目のは$-30$。}であれば是非トップになってほしい。まずは門前でチートイ以外はほぼ無理。\mj{9p} の対子と役牌を落とすと僅かな確率でタンヤオがつく。役牌で和了れば、浮き牌の \mj{2s} とドラ表示牌の \mj{7z} が打牌候補になる\citep[][pp. 114--119]{tetsusenryaku}。

\onlyinsubfile{\bibliographystyle{../jecon}%
\nobibliography{../refmj}}


\end{document}